\chapter{Реализация консольного приложения} \label{ch5}

В данной главе описана реализация консольного приложения, разработанного по принципам TDD.
Тесты, спроектированные в главе~\ref{ch2} и реализованные в главе~\ref{ch4}, выступают как исполняемая спецификация: каждый модуль создаётся ровно в том объёме, который необходим для прохождения соответствующих тест-кейсов.
Проект реализуется на JavaScript в среде Node.js; запуск тестов осуществляется командой \texttt{npm test}, а запуск приложения -- командой \texttt{node src/index.js}.

\section{Структура проекта} \label{ch5:sec1}

Исходный код распределён по нескольким каталогам, каждый из которых отвечает за строго определённый слой приложения.
Тестовый код зеркалирует структуру основного кода: тестовые файлы располагаются в каталоге \texttt{tests} и сгруппированы по назначению.
Структура проекта представлена в \taref{tab:packages}.

\noindent
\begingroup
\centering
\small
\begin{longtable}[c]{|p{7.8cm}|p{8.2cm}|}
    \caption{Структура проекта}
    \label{tab:packages}
    \\
    \hline
    Каталог / модуль & Назначение \\ \hline
    \endfirsthead
    \captionsetup{format=tablenocaption,labelformat=continued}
    \caption[]{}\\
    \hline
    Каталог / модуль & Назначение \\ \hline
    \endhead
    \hline
    \endfoot
    \hline
    \endlastfoot
    \texttt{src/index.js} & Точка входа в приложение \\ \hline
    \texttt{src/ui/consoleUi.js} & Главное меню и обработка пользовательского ввода \\ \hline
    \texttt{src/analyzers/manacherOdd.js} & Вычисление массива нечётных радиусов \\ \hline
    \texttt{src/analyzers/manacherEven.js} & Вычисление массива чётных радиусов \\ \hline
    \texttt{src/analyzers/centerExpansion.js} & Эталонный алгоритм расширения от центра \\ \hline
    \texttt{src/services/palindromeService.js} & Формирование агрегированного результата анализа \\ \hline
    \texttt{src/io/jsonReader.js} & Загрузка входной строки из JSON-файла \\ \hline
    \texttt{src/io/resultWriter.js} & Вывод результата в консоль и запись результата в JSON-файл \\ \hline
    \texttt{tests/*.test.js} & Набор модульных тестов \\ \hline
\end{longtable}
\normalsize
\endgroup

\section{Описание основных модулей} \label{ch5:sec2}

Назначение ключевых модулей приложения приведено в \taref{tab:classes}.

\noindent
\begingroup
\centering
\small
\begin{longtable}[c]{|p{5.5cm}|p{2cm}|p{8cm}|}
    \caption{Основные модули приложения и их назначение}
    \label{tab:classes}
    \\
    \hline
    Модуль & Тип & Назначение \\ \hline
    \endfirsthead
    \captionsetup{format=tablenocaption,labelformat=continued}
    \caption[]{}\\
    \hline
    Модуль & Тип & Назначение \\ \hline
    \endhead
    \hline
    \endfoot
    \hline
    \endlastfoot
    \texttt{index.js} & Модуль & Создаёт экземпляр интерфейса и запускает приложение \\ \hline
    \texttt{consoleUi.js} & Модуль & Отображает меню, считывает команды пользователя, управляет сценарием работы \\ \hline
    \texttt{manacherOdd.js} & Модуль & Вычисляет массив радиусов нечётных палиндромов \\ \hline
    \texttt{manacherEven.js} & Модуль & Вычисляет массив радиусов чётных палиндромов \\ \hline
    \texttt{centerExpansion.js} & Модуль & Реализует базовый конкурирующий алгоритм расширения от центра \\ \hline
    \texttt{palindromeService.js} & Модуль & Объединяет результаты алгоритмов, восстанавливает максимум и подсчитывает статистику \\ \hline
    \texttt{jsonReader.js} & Модуль & Потоково читает и валидирует входной JSON-файл \\ \hline
    \texttt{resultWriter.js} & Модуль & Печатает результат анализа в консоль и сериализует его в JSON \\ \hline
\end{longtable}
\normalsize
\endgroup

Разделение проекта на небольшие модули уменьшает связанность кода и упрощает тестирование.
Каждый модуль отвечает за ограниченный набор обязанностей, что особенно важно при применении TDD.

\section{Описание реализации вычислительной части} \label{ch5:sec3}

Вычислительное ядро проекта состоит из трёх анализаторов:
\begin{enumerate}
    \item модуль расширения от центра;
    \item модуль вычисления нечётных радиусов алгоритмом Манакера;
    \item модуль вычисления чётных радиусов алгоритмом Манакера.
\end{enumerate}

Базовый алгоритм расширения от центра используется как эталон при тестировании.
Он проходит по каждому возможному центру и последовательно расширяет границы палиндрома.
Хотя этот подход менее эффективен, он удобен как проверочная реализация.

Основной производственный алгоритм реализован двумя модулями: \texttt{manacherOdd.js} и \texttt{manacherEven.js}.
Первый рассчитывает массив $d_{odd}$, второй -- массив $d_{even}$.
Это разделение упрощает код и делает границы ответственности прозрачными.

Сервисный модуль \texttt{palindromeService.js} выполняет:
\begin{enumerate}
    \item запуск обоих модулей Манакера;
    \item запуск эталонного алгоритма расширения от центра при необходимости сравнения;
    \item поиск наибольшего палиндрома;
    \item подсчёт общего числа палиндромных подстрок;
    \item формирование рекомендаций пользователю;
    \item подготовку итогового объекта результата.
\end{enumerate}

\section{Тестирование} \label{ch5:sec4}

Тестовый набор включает модульные тесты для всех ключевых частей приложения:
\begin{itemize}
    \item вычисления нечётных радиусов;
    \item вычисления чётных радиусов;
    \item эталонного алгоритма расширения от центра;
    \item сервисного слоя;
    \item чтения JSON-файла;
    \item записи результата в файл и вывода в консоль;
    \item пользовательского интерфейса.
\end{itemize}

Запуск тестов осуществляется командой \texttt{npm test}.
Тесты проверяют не только корректность результата, но и совпадение ответов алгоритма Манакера с ответами эталонного алгоритма на тех же строках.

\section{Демонстрация работы приложения} \label{ch5:sec5}

Приложение запускается командой \texttt{node src/index.js}.
После запуска отображается приветственное сообщение и главное меню.
В верхней части меню выводится информация о текущем состоянии: введена ли строка, какова её длина и доступен ли последний рассчитанный результат.

Главное меню приложения представлено в \taref{tab:runtime-menu}.

\noindent
\begingroup
\centering
\small
\begin{longtable}[c]{|p{2.2cm}|p{13cm}|}
    \caption{Пример структуры отображаемого меню}
    \label{tab:runtime-menu}
    \\
    \hline
    Элемент & Содержание \\ \hline
    \endfirsthead
    \captionsetup{format=tablenocaption,labelformat=continued}
    \caption[]{}\\
    \hline
    Элемент & Содержание \\ \hline
    \endhead
    \hline
    \endfoot
    \hline
    \endlastfoot
    Заголовок & Название приложения и краткое описание задачи \\ \hline
    Статус & Строка не загружена / загружена, длина строки, наличие последнего результата \\ \hline
    Пункты меню & Ввод вручную, загрузка JSON, запуск Манакера, запуск эталонного алгоритма, вывод отчёта, сохранение результата, справка, выход \\ \hline
\end{longtable}
\normalsize
\endgroup

При выборе пункта ручного ввода пользователь вводит строку напрямую.
При выборе загрузки из файла приложение ожидает путь к JSON-файлу, после чего выполняет его потоковое чтение и валидацию.
Если файл некорректен, пользователю выводится диагностическое сообщение без аварийного завершения программы.

После выбора алгоритма приложение формирует результат анализа.
В консоль выводятся:
\begin{itemize}
    \item исходная строка;
    \item массивы $d_{odd}$ и $d_{even}$;
    \item наибольшая палиндромная подстрока;
    \item её координаты;
    \item суммарное число палиндромных подстрок;
    \item рекомендация пользователю.
\end{itemize}

На \firef{fig:result-flow} приведена блок-схема формирования итогового результата в модуле сервиса.

\begin{figure}[h]
    \centering
    \begin{tikzpicture}[
        node distance=6mm and 8mm,
        block/.style={rectangle, draw, rounded corners, align=center, minimum width=4.6cm, minimum height=0.9cm},
        decision/.style={diamond, draw, aspect=2, align=center, inner sep=1pt},
        line/.style={-Latex}
    ]
        \node[block] (input) {Получить строку, $d_{odd}$ и $d_{even}$};
        \node[block, below=of input] (search) {Найти максимальный радиус\\и тип палиндрома};
        \node[decision, below=of search] (type) {Палиндром\\чётный?};
        \node[block, below left=of type] (odd) {Восстановить границы\\нечётного палиндрома};
        \node[block, below right=of type] (even) {Восстановить границы\\чётного палиндрома};
        \node[block, below=16mm of type] (count) {Подсчитать суммарное\\число палиндромов};
        \node[block, below=of count] (result) {Сформировать объект\\итогового результата};

        \draw[line] (input) -- (search);
        \draw[line] (search) -- (type);
        \draw[line] (type) -- node[above left]{нет} (odd);
        \draw[line] (type) -- node[above right]{да} (even);
        \draw[line] (odd) |- (count);
        \draw[line] (even) |- (count);
        \draw[line] (count) -- (result);
    \end{tikzpicture}
    \caption{Блок-схема формирования итогового результата анализа строки}
    \label{fig:result-flow}
\end{figure}

\FloatBarrier

Пункт сохранения результата записывает последний рассчитанный анализ в JSON-файл.
В сохраняемой структуре присутствуют поля \texttt{algorithm}, \texttt{oddRadii}, \texttt{evenRadii}, \texttt{longestPalindrome}, \texttt{palindromeCount}, \texttt{recommendation} и, для основного алгоритма, объект \texttt{comparison}.
Если результат ещё не сформирован, приложение выводит сообщение о необходимости сначала выполнить анализ.

Пункт справки содержит:
\begin{itemize}
    \item описание входного формата JSON;
    \item пояснение массивов радиусов;
    \item описание рекомендации;
    \item пример структуры результирующего файла.
\end{itemize}

\section{Выводы} \label{ch5:conclusion}

В данной главе описана реализация консольного приложения на JavaScript, разработанного по принципам TDD.
Проект разделён на независимые модули: анализаторы, сервисный слой, ввод-вывод и пользовательский интерфейс.
В приложении реализованы два способа решения задачи -- алгоритм Манакера и расширение от центра, что соответствует требованиям задания и упрощает тестирование через сравнение результатов.
Подробные листинги основного кода и тестов приведены в приложениях.
